\documentclass{article}
\usepackage{amsmath}
\usepackage{graphicx}
\usepackage{geometry}
\geometry{a4paper, margin=1in}

\title{MORGN: Momentum Orthogonalization via Recursive Gauss-Newton}
\author{Your Name}

\begin{document}

\maketitle

\begin{abstract}
We present MORGN, a second-order optimization algorithm for training deep neural networks. MORGN is a quasi-Newton method that approximates the inverse Hessian matrix using a recursive Gauss-Newton (RLS-style) update. This allows the optimizer to take more effective, Newton-like steps, which can lead to faster convergence and better performance on difficult optimization problems.
\end{abstract}

\section{Introduction}

Second-order optimization methods have the potential to significantly accelerate the training of deep neural networks. However, traditional second-order methods, such as Newton's method, are often too computationally expensive to be practical. Quasi-Newton methods, which approximate the Hessian matrix instead of computing it directly, offer a more computationally efficient alternative.

MORGN is a quasi-Newton method that uses a recursive Gauss-Newton (RLS-style) update to approximate the inverse Hessian matrix. This allows the optimizer to build up a more accurate approximation of the inverse Hessian over time, while still remaining computationally efficient.

\section{Algorithm}

The core of the MORGN optimizer is the left-preconditioned update for 2D parameters. For each matrix parameter $W \in \mathbb{R}^{m \times n}$, the optimizer maintains an inverse Hessian approximation $P \in \mathbb{R}^{m \times m}$ on the smaller/left dimension (transposing if $m > n$). The update is then computed as follows:

\begin{enumerate}
    \item Compute the preconditioned step: $U = P G$, where $G = \frac{\partial L}{\partial W}$ is the gradient of the loss with respect to $W$.
    \item Update the parameter: $W \leftarrow W - \eta U$, where $\eta$ is the learning rate.
\end{enumerate}

The inverse Hessian approximation $P$ is updated at each step using a rank-1 RLS/Sherman-Morrison formula. This update incorporates new gradient information and allows the optimizer to build up a more accurate approximation of the inverse Hessian over time.

\section{Conclusion}

MORGN is a promising new optimization algorithm for training deep neural networks. It is a computationally efficient quasi-Newton method that has the potential to significantly accelerate the training of deep neural networks. Further research is needed to evaluate the performance of MORGN on a wide range of deep learning tasks.

\end{document}
